\section{Failure Model and Specification}
We follow \texttt{FAILURES\_SPEC.md} v0.2.0, the single source of truth.

\paragraph{Failure taxonomy.} Hard failure (500/exception), \emph{ambiguous outcome} (caller does not know if side effect happened---timeout, crash before ack, partition), partial success, duplicate execution (retry/redelivery).

\paragraph{Principle.} A named, testable invariant: \{id, question, invariant, applies\_to, failure\_modes, patterns, tests\}. Example---idempotency: ``What happens if this operation executes more than once?'' Invariant: one logical operation yields at most one effect.

\paragraph{11 dimensions.} atomicity, idempotency, ordering, concurrency, availability, timeout, consistency, resource\_exhaustion, recovery, observability, retry\_safety. Each defines a builder question (e.g., timeout: ``What happens when you don't know whether it succeeded?'').

\paragraph{Finding.} Unit of output: \{id, dimension, severity (CRITICAL/HIGH/MEDIUM/LOW), confidence 0--1, mode, title, why, evidence\{excerpt, lines\}, risk, required/optional mitigations, tests\}. Severity is impact-anchored; confidence is honest.

\paragraph{Invariants.} Predicates that must hold under \emph{any} failure, e.g., \texttt{invariant.payment.not\_double\_charged}, \texttt{invariant.enrollment.requires\_successful\_payment}. \texttt{check\_invariant} returns violatable? + counter-scenario.
